\chapter{Denoising and Matrix Completion}
\label{ch:denoising-matrix-completion}

\section{Optional Module}

This chapter is optional. It uses the SVD as a practical tool, but it is not a
new theory unit. The goal is to reinforce one idea:
\[
  \text{large singular values often carry structure; small ones often carry noise.}
\]
That sentence is not a theorem in every situation. It is a modeling principle
that works when the signal is approximately low rank.

\section{Low-Rank Signal Plus Noise}

Suppose an observed matrix has the form
\[
  M=A+E,
\]
where $A$ is a low-rank signal and $E$ is noise. If $A$ has rank $r$, then its
SVD has only $r$ nonzero singular values. Noise usually spreads energy across
many directions.

In a favorable case, the singular values of $M$ look like
\[
  \sigma_1\geq\cdots\geq\sigma_r
  \quad\text{large, then a drop, then a tail.}
\]
The drop is often called an \emph{elbow}.

\section{Truncated SVD Denoising}

The denoising procedure is:
\begin{enumerate}
  \item compute the SVD $M=U\Sigma V^T$;
  \item choose a rank $k$;
  \item reconstruct
        \[
          M_k=\sum_{i=1}^k \sigma_i u_i v_i^T.
        \]
\end{enumerate}

If $k$ is too small, important signal is removed. If $k$ is too large, noise is
kept.

\begin{warningbox}
SVD truncation is not automatic truth extraction. It is a good method when the
signal is well approximated by a low-rank matrix and the noise is not aligned
with the leading singular directions.
\end{warningbox}

\section{A Basic NumPy Experiment}

\begin{lstlisting}
import numpy as np

rng = np.random.default_rng(0)

m, n, r = 40, 30, 2
L = rng.standard_normal((m, r))
R = rng.standard_normal((r, n))
A_true = L @ R                       # rank at most 2

noise = 0.5 * rng.standard_normal((m, n))
M = A_true + noise

U, s, Vt = np.linalg.svd(M, full_matrices=False)
print("first singular values:", s[:8])

for k in [1, 2, 5, 10]:
    M_k = U[:, :k] @ np.diag(s[:k]) @ Vt[:k, :]
    err_to_signal = np.linalg.norm(M_k - A_true, "fro")
    err_to_observed = np.linalg.norm(M_k - M, "fro")
    print(k, err_to_signal, err_to_observed)
\end{lstlisting}

The best $k$ for recovering the original signal need not be the best $k$ for
approximating the noisy observation. Approximating the observation too closely
can mean fitting noise.

\section{Choosing The Rank}

There is no universal rank-choice rule for an introductory course. Three
reasonable first checks are:

\begin{itemize}
  \item look for an elbow in the singular values;
  \item keep enough singular values to explain a chosen energy fraction;
  \item test performance on held-out entries or held-out data.
\end{itemize}

The energy fraction of the top $k$ singular values is
\[
  \frac{\sigma_1^2+\cdots+\sigma_k^2}
       {\sigma_1^2+\cdots+\sigma_r^2}.
\]

\section{Matrix Completion}

Matrix completion asks for missing entries of a matrix. Without assumptions,
this is impossible: missing numbers could be anything. Low-rank structure is
one useful assumption.

For example, suppose a user-item rating matrix is approximately low rank
because a few hidden preference factors explain many ratings. Then the observed
entries may contain enough information to estimate the missing entries.

At this level, the linear algebra message is modest:
\[
  \text{low rank reduces degrees of freedom.}
\]
A general $m\times n$ matrix has $mn$ entries. A rank-$k$ factorization
\[
  A=LR,
  \qquad
  L\in\R^{m\times k},
  \quad
  R\in\R^{k\times n},
\]
uses about $k(m+n)$ parameters. When $k\ll \min(m,n)$, that is much smaller.

\section{A Simple Completion Objective}

Let $\Omega$ be the set of observed entries. One basic low-rank fitting problem
is
\[
  \min_{L,R}
  \sum_{(i,j)\in\Omega}
  \left((LR)_{ij}-M_{ij}\right)^2.
\]
This is not a linear least-squares problem in $L$ and $R$ together, because the
unknowns multiply each other. But if $R$ is fixed, solving for $L$ is least
squares; if $L$ is fixed, solving for $R$ is least squares. More advanced
matrix completion algorithms build on this idea.

\begin{aimlbox}
Modern low-rank methods in machine learning often use the same parameter-count
idea: represent a large update by two thin matrices instead of one full matrix.
That is an application of rank intuition, not a replacement for learning the
linear algebra itself.
\end{aimlbox}

\begin{labbox}
Lab 12 uses synthetic low-rank matrices plus noise and compares SVD
truncations. The experiment is intentionally small enough that every step can
be inspected with basic NumPy.
\end{labbox}

\section*{Chapter Summary}
\addcontentsline{toc}{section}{Chapter Summary}

\begin{itemize}
  \item Low-rank matrices have only a few nonzero singular values.
  \item Noise often spreads across many singular directions.
  \item Truncated SVD can denoise when signal and noise separate spectrally.
  \item Choosing the truncation rank is a modeling decision.
  \item Matrix completion needs assumptions; low rank is one useful assumption.
\end{itemize}

\section*{Exercises}
\addcontentsline{toc}{section}{Exercises}

\subsection*{A. Theory}

\begin{exercise}
\exerciseid{16.A1}
Explain why a low-rank signal plus full-rank noise may still have a visible
singular-value elbow.
\end{exercise}

\begin{exercise}
\exerciseid{16.A2}
Explain why a rank-$k$ factorization $LR$ uses far fewer parameters than a full
$m\times n$ matrix when $k\ll \min(m,n)$.
\end{exercise}

\subsection*{B. By Hand}

\begin{exercise}
\exerciseid{16.B1}
Suppose a matrix has singular values $10,7,1,0.8,0.7$. Which rank would you try
first for denoising, and why?
\end{exercise}

\begin{exercise}
\exerciseid{16.B2}
A $100\times 80$ rank-$3$ factorization $LR$ uses
$L\in\R^{100\times 3}$ and $R\in\R^{3\times 80}$. Compare the number of
parameters with the full $100\times 80$ matrix.
\end{exercise}

\subsection*{C. Python / NumPy}

\begin{exercise}
\exerciseid{16.C1}
Generate a rank-$2$ matrix plus Gaussian noise. Compute singular values and
compare rank-$1$, rank-$2$, and rank-$5$ reconstructions.
\end{exercise}

\begin{exercise}
\exerciseid{16.C2}
For your noisy matrix from 16.C1, compute the energy fraction kept by the top
$k$ singular values for $k=1,2,5$.
\end{exercise}

\subsection*{D. Challenge}

\begin{exercise}
\exerciseid{16.D1}
Explain why too small a truncation rank underfits while too large a rank leaves
noise in the reconstruction.
\end{exercise}

